\chapter{Conclusions}

The custom set-up for the measurement of dynamic impedance is very expensive. Miniaturization via integration to reduce cost is possible but require the investment of companies or reserach centers. The investment is justified when the information obtain by the low frequency band give enough advanced understanding on the kinetics of the systems. Batteries are those systems where the low frequency band is of interest and might drive the development of this technique. The prototypation stage of this work allow a convenient acquisition of the impedance over a broad band of frequency lowering the entering barrier for electrochesmists to apprach the method. When measruing impedance will become more present in the literature, for application even behiond my immagination, will be the real word test of the potentiality of this method. As nowaday no electrochemsitry thinks of the inner mechanisms of an electrochemical workstation or even the computer that controls it or analyzez the data, the same fate shall be for the measurement of non-stationary impedance.
While the community of system engeneers have develpooed important tools for the identifaction of time varying trasnfer function, the knowlkedge is still segregated from the electrochemistry community for its combersom language and specific formalisms. The scintific community needs a link of communication between this to words to allow a real impact on the reserach. Usually this happen when a commercial device automate the usage of the method. On the other hand no investments will be carries when the experimental electrochemsitry have no understanding or interest on the methodology. This is the ground this thesis is trying to place, surly failing in the system identification realm npot having myself the background knosledge to understand the formalisms of functional analysis.